% SPDX-License-Identifier: GPL-3.0-or-later
% Copyright (C) 2026 Alessandro Greco
\section*{Глава \ref{cap:19-rete-privata}: приватная сеть Tor}

\begin{soluzione}{es:19-porte}
\texttt{test006e} — узел номер 6: ORPort 5106, ControlPort 8006. У \texttt{test009h} нет ORPort: onion-сервис не является ретранслятором и не принимает соединения от других узлов сети. Он сам строит исходящие цепи к точкам представления и встречи.
\end{soluzione}

\begin{soluzione}{es:19-guard}
Это параметр \opt{TestingDirAuthVoteGuard}~\texttt{*} в файлах torrc центров каталогов. В публичной сети флаг Guard требует измеренных за некоторое время стабильности и пропускной способности (\cref{cap:12-directory}). В лаборатории любой ретранслятор может быть выбран входным узлом: статистика выбора первого перехода не соответствует реальной сети.
\end{soluzione}

\begin{soluzione}{es:19-socks}
При \texttt{--socks5} программа \texttt{curl} самостоятельно разрешает имя \texttt{localhost} и передаёт клиенту Tor IP-адрес. При \texttt{--socks5-hostname} передаётся имя, которое разрешает выходной ретранслятор. При посещении обычных сайтов в публичной сети первый способ выдавал бы доменные имена локальному DNS-резолверу и его наблюдателям, хотя сам трафик шёл бы через Tor. Это классическая утечка DNS.
\end{soluzione}

\begin{soluzione}{es:19-consenso}
Каждые 20 секунд согласно \opt{V3AuthVotingInterval}. За сутки это $86400/20 = 4320$ консенсусов в лаборатории против 24 в публичной сети, где интервал голосования равен одному часу.
\end{soluzione}

\begin{soluzione}{es:19-autorita}
При остановке одного центра каталогов сеть продолжает выпускать консенсусы, поскольку достаточно большинства подписей (две из трёх). При остановке двух единственный оставшийся центр не достигает большинства, и новый действительный консенсус не создаётся. Клиенты, однако, продолжают использовать старый консенсус как \emph{reasonably live} — ещё примерно 24 часа после истечения его срока (\cref{cap:12-directory}) — и могут долго строить цепи. В лаборатории срок консенсуса составляет 40 секунд, но допустимое использование после окончания срока остаётся равным 24 часам.
\end{soluzione}

\begin{soluzione}{es:19-stem}
План: \file{c.get\_network\_status(fingerprint)} возвращает запись консенсуса с полями \texttt{address} и \file{or\_port}. Вызовите метод для каждого отпечатка из \texttt{circ.path}, то есть для первого элемента каждой пары.
\end{soluzione}

\begin{soluzione}{es:19-rete}
Определение сети сохраняет структуру из книги, но использует \texttt{getN(5)}, \texttt{getN(10)} и \texttt{getN(4)}. Время начальной загрузки зависит от компьютера. При увеличении числа узлов особенно растёт ожидание появления ретрансляторов в консенсусе и загрузки микродескрипторов клиентами.
\end{soluzione}

\begin{soluzione}{es:19-log}
Исследовательское упражнение: найдите сообщения о выборе маршрута, отправке \cella{CREATE2}/\texttt{EXTEND2}, получении ответов и завершении построения цепи. Уровень журналирования \emph{info} показывает не все ячейки; более подробные сведения можно получить, включив \emph{debug} в torrc клиента.
\end{soluzione}

\begin{soluzione}{es:19-cattura}
Нет: после рукопожатия TLS ячейки передаются зашифрованными, а \cmd{tcpdump} видит только записи TLS. При этом остаются видимыми адреса, порты, размеры записей и время — данные, на которых основаны атаки анализа трафика (\cref{cap:26-attacchi}).
\end{soluzione}

\section*{Глава \ref{cap:21-onion-pratica}: публикация onion-сервиса}

\begin{soluzione}{es:21-vanity}
Шесть символов заданного префикса требуют в среднем $32^6 \approx 1{,}07 \times 10^9$ попыток, около 107 секунд при десяти миллионах попыток в секунду. Восемь символов требуют $32^8 \approx 1{,}1 \times 10^{12}$ попыток, около $1{,}1 \times 10^5$ секунд, чуть более 30 часов. Каждый дополнительный символ увеличивает время в 32 раза.
\end{soluzione}

\begin{soluzione}{es:21-fuga}
Страница ошибки раскрывает программное обеспечение и версию веб-сервера (это помогает искать уязвимости) и, что ещё важнее, IP-адрес 192.0.2.10, который может оказаться настоящим публичным адресом сервера: тогда сервис можно локализовать. Необходимо скрыть сигнатуры сервера в заголовках и страницах ошибок, разрешить ему слушать только localhost или Unix-сокет и проверить ответы лабораторным клиентом.
\end{soluzione}

\begin{soluzione}{es:21-revoca}
Удалите \file{alice.auth} и перечитайте конфигурацию. Сервис публикует новый дескриптор без записи Алисы, зашифрованный новым \emph{descriptor cookie}, который она не может расшифровать. Клиенты со старым дескриптором могут пользоваться известными точками представления, пока они действуют; для более полного отзыва стоит сменить точки представления, например перезапустив сервис. Крайний случай: если Алиса была единственным авторизованным клиентом, каталог окажется пустым, а HSDir отклонят дескриптор (\cref{sec:onionp-ristretta}). Для настоящего отключения ограниченного обнаружения удалите каталог \file{authorized\_clients/} и перечитайте конфигурацию.
\end{soluzione}

\begin{soluzione}{es:21-socket}
Задайте \texttt{HiddenServicePort 80 unix:/path/socket} и настройте веб-сервер на прослушивание этого Unix-сокета. Помимо отсутствия открытого TCP-порта, сокет защищается правами файловой системы, недоступен через сетевой интерфейс и исключает случайное открытие сервиса на публичном интерфейсе.
\end{soluzione}

\section*{Глава \ref{cap:20-traffico}: наблюдение за трафиком}

\begin{soluzione}{es:20-dimensioni}
TLS не выравнивает границы записей по ячейкам Tor. tor помещает ячейки в буфер, а библиотека TLS объединяет их в записи длиной до 16\,KB (отсюда значение \texttt{sendme\_inc}, равное 31 ячейке, \cref{sec:cc}) и добавляет заголовок и тег аутентификации. TCP может делить или объединять эти записи в сегментах. Фиксированная длина относится к ячейкам, а не к сетевым пакетам.
\end{soluzione}

\begin{soluzione}{es:20-relay}
План: после \cella{CREATED2} расшифруйте ячейки \cella{RELAY\_EARLY} и \cella{RELAY} лабораторной цепи, снимая слои шифрования по порядку с помощью ключей переходов из \file{circuit.py}. Состояние потока AES-CTR должно продвигаться ровно в порядке использования. При построении цепи видны \texttt{EXTEND2} и \texttt{EXTENDED2}, затем возможные \texttt{BEGIN}, \texttt{CONNECTED}, \texttt{DATA} и \texttt{SENDME}.
\end{soluzione}

\begin{soluzione}{es:20-padding}
На клиентском канале, по которому проходят цепи из нескольких переходов, записи TLS могут появляться через несколько секунд даже без прикладного трафика. Это ячейки заполнения канала \cella{PADDING}, передаваемые через случайные интервалы от 1,5 до 9,5 секунды (\cref{sec:lab-cifratura}). На канале, не использовавшемся для многошаговых цепей, такого трафика нет.
\end{soluzione}

\begin{soluzione}{es:20-dissector}
Задание по программированию: анализатор Lua, зарегистрированный как \emph{эвристический диссектор} расшифрованных данных TLS, читает первые два байта CircID начальной ячейки \cella{VERSIONS} или четыре байта в последующих ячейках, затем байт команды и поле длины переменных ячеек. Нужно сохранять состояние отдельно для каждого соединения и направления, как в \file{traffic.py}.
\end{soluzione}

\section*{Глава \ref{cap:22-relay}: управление ретранслятором или мостом}

\begin{soluzione}{es:22-ordine}
Да. Первое подходящее правило \texttt{accept *:80} разрешает соединение, последующие правила уже не проверяются. Достаточно изменить порядок: \texttt{reject 192.0.2.0/24:*}, \texttt{accept *:80}, \texttt{reject *:*}.
\end{soluzione}

\begin{soluzione}{es:22-riassunto}
Сводка содержит только порты, без адресов, и является приближённой. Порт считается разрешённым для \enquote{большинства} адресов согласно точному правилу dir-spec, а не просто при доле свыше 50\%: выходной ретранслятор не должен блокировать более $2^{25}$ IPv4-адресов (одна сеть /7) или $2^{112}$ IPv6-адресов (одна сеть /16), не считая непубличных адресов. Получив \texttt{END} с причиной \texttt{EXITPOLICY}, клиент на время TTL запоминает отказ выходного ретранслятора для указанного адреса и порта и повторяет запрос в другой цепи через другой выходной узел.
\end{soluzione}

\begin{soluzione}{es:22-famiglia}
При \opt{MyFamily} учитываются только взаимные объявления: пять прежних ретрансляторов принадлежат к одному семейству, но шестой — нет, поскольку ни один из прежних его не указал. Поэтому клиент может выбрать шестой узел вместе с другим членом семейства в одной цепи. С ключом семейства достаточно установить этот ключ на новом ретрансляторе: подписанный сертификат связывает его с семейством, не требуя изменений на остальных узлах.
\end{soluzione}

\begin{soluzione}{es:22-refused}
В журналах последовательно появляются \enquote{in-progress connect failed \ldots (Connection refused)} и \enquote{Didn't recognize errno 115}. Если причину окончания потока вычислять по ошибке, которую возвращает \texttt{getsockopt(SO\_ERROR)}, клиент получит \texttt{CONNECTREFUSED}. Аналогично можно верно передавать \texttt{TIMEOUT} при \texttt{ETIMEDOUT}, \texttt{NOROUTE} при \texttt{ENETUNREACH} или \texttt{EHOSTUNREACH} и \texttt{CONNRESET}, вместо \texttt{MISC} для соединений, завершившихся ошибкой после \texttt{EINPROGRESS}. Для клиента меняется сообщение об ошибке, а не сам порядок обработки неудачи.
\end{soluzione}

\begin{soluzione}{es:22-accounting}
При 10\,MB/s ретранслятор передаёт около 864\,GB в сутки; учитывается максимум между отправленным и полученным трафиком. Лимит 1\,TB исчерпывается чуть более чем за день. Ретранслятор, выключенный большую часть месяца, теряет флаги \texttt{Stable} и \texttt{Guard} и разрывает цепи при остановке; медленный узел, непрерывно работающий весь месяц, гораздо полезнее. Поэтому \opt{RelayBandwidthRate} стоит настроить так, чтобы лимита хватило на весь период.
\end{soluzione}
